% 已完成自动评测；语义编辑核验仍待完成。
\section{实验}
我们在三个DL3DV室内开发场景上比较基础单Agent、工具单Agent和完整WorldHarness，采用Astra与Sol 6.1两种建模模型。主对照采用ultra推理强度，完整框架另设high消融，共24个条件。每条件保存M0、M1、M2三版；完整框架执行R1、R2两次独立审核，单Agent自行返修，不接收系统生成的叠图或外部审核意见。返修与审核没有硬性时间上限。

本批已完成全部72份自动评测。每场景48张构建图、6张可见验证图，另用12个私有测试视角；测试统一为640像素宽、32次Cycles采样和相同OIDN。报告PSNR、SSIM和AlexNet LPIPS。空间指标采用留出图像三视角三角化的稀疏参考：首交点相对误差度量射线第一表面的距离偏差，点到表面相对距离度量最近表面贴合；均按参考射线长度归一化，截断为1，缺失记1，并按帧等权平均。该参考不提供完整扫描几何真值。

\paragraph{执行历史与限制。}
18个M0继承经哈希验证的v1首版，其他条件续接或补齐，因此初始预算具有历史差异。已接触外部反馈的单Agent旧返修被隔离，从干净M0会话继续。执行期间放宽了报告的计划观测绑定，并将离线GLB上限与工作台导入上限分开；原失败、恢复位置和消耗均保留。各条件仅一次运行，且完整框架同时改变计划、初始化和审核，结果不能作为方差受控或单组件归因的最终结论。

\begin{table*}[t]
\centering
\caption{M2的三场景等权平均。编辑率为机械操作通过率，尚不代表目标语义归属正确。}
\begin{tabular}{lllrrrrrr}
\hline
模型 & 条件 & 强度 & PSNR$\uparrow$ & SSIM$\uparrow$ & LPIPS$\downarrow$ & 首交点$\downarrow$ & 表面距离$\downarrow$ & 编辑$\uparrow$ \\
\hline
Astra & 基础单Agent & ultra & 19.10 & 0.7486 & 0.2481 & 0.0567 & 0.0161 & 100.0\% \\
Astra & 工具单Agent & ultra & 19.56 & 0.7661 & 0.2333 & 0.0432 & 0.0107 & 100.0\% \\
Astra & 完整Harness & ultra & 13.67 & 0.6420 & 0.3580 & 0.0496 & 0.0126 & 93.3\% \\
Sol 6.1 & 基础单Agent & ultra & 18.22 & 0.7190 & 0.2730 & 0.0467 & 0.0091 & 100.0\% \\
Sol 6.1 & 工具单Agent & ultra & 19.00 & 0.7627 & 0.2493 & 0.0425 & 0.0116 & 100.0\% \\
Sol 6.1 & 完整Harness & ultra & 14.03 & 0.6536 & 0.3541 & 0.0472 & 0.0114 & 93.3\% \\
Astra & 完整Harness & high & 14.88 & 0.6683 & 0.3405 & 0.0469 & 0.0107 & 93.3\% \\
Sol 6.1 & 完整Harness & high & 12.25 & 0.4824 & 0.4625 & 0.0393 & 0.0108 & 100.0\% \\
\hline
\end{tabular}
\end{table*}

\paragraph{结果分析。}
本批当前实现没有显示出独立审核返修对外观的稳定收益。几何误差存在改善与退步，首交点与最近表面距离需同时报告，覆盖率也不能单独证明空间结构正确。当前只完成机械编辑的移动、旋转、缩放、隐藏、改材质及保存重开检查；目标语义归属和几何拆分完整性待独立视觉核验，物理真实性不在本轮评价范围内。

\paragraph{消耗。}
输入、输出、缓存输入、推理输出和实际调用耗时按角色记录。Token采用每个模型线程的最后累计快照去重，缓存与推理为子项，不重复相加。各条件的调用累计耗时不同于并行批次墙钟时间；共享初始化与隔离诊断成本单列。
